\chapter{Introduction}
Distributed consensus is a foundational problem in distributed computing: a set
of nodes must agree on a single value even when some nodes may fail or messages
may be lost or delayed. Classical results such as the FLP impossibility theorem
show that deterministic consensus cannot be guaranteed in a fully asynchronous
system with even a single faulty process \cite{Fischer1985}. Practical systems
work around this by making partial synchrony assumptions, but these assumptions
are frequently violated in real wide-area deployments.

\section{Motivation}
Modern distributed applications increasingly run across wide-area networks with
variable latency and non-trivial packet loss. Protocols designed and tuned for
data-centre networks, where loss is rare and latency is predictable, often behave
poorly once deployed across such links \cite{Ongaro2014}.

\section{Contributions}
This thesis makes the following contributions:
\begin{itemize}
  \item A formal model of consensus under bounded but non-zero message loss.
  \item A protocol whose expected termination time degrades smoothly with loss
    rate, rather than failing once a fixed threshold is crossed.
  \item An experimental evaluation comparing the protocol against Raft-style
    baselines \cite{Ongaro2014} under simulated partition and loss.
\end{itemize}

% A figure would normally be included here, e.g.:
% \begin{figure}[htbp]
%   \centering
%   \includegraphics[width=0.7\textwidth]{figures/architecture}
%   \caption{System architecture overview.}
%   \label{fig:architecture}
% \end{figure}

\begin{table}[htbp]
  \centering
  \caption{Notation used throughout this thesis.}
  \label{tab:notation}
  \begin{tabular}{ll}
    \toprule
    Symbol & Meaning \\
    \midrule
    $n$ & number of nodes \\
    $f$ & number of tolerated faulty nodes \\
    $p$ & message loss probability \\
    \bottomrule
  \end{tabular}
\end{table}

% An algorithm listing would normally be included here, e.g. using the
% algorithm2e or algorithmicx package.

